\section{Jerusalem, Judea, and the nations}\label{sec:peter-acts}

Acts makes Peter the principal apostolic speaker in its first half. This is a
literary and theological presentation by Luke, not a stenographic archive of
every speech. Yet its broad picture---Peter publicly proclaiming Jesus,
sharing leadership with the Twelve, encountering both Jewish and Gentile
hearers, and eventually yielding the narrative foreground to Paul---is
consistent with the independent fact that Paul sought Cephas in Jerusalem
and later counted him among the pillars.

\subsection{The Jerusalem witness}

Before Pentecost Peter interprets Judas's death and directs the choice of
Matthias to restore the Twelve. At Pentecost he stands with the Eleven and
proclaims that the crucified and risen Jesus is Lord and Messiah. After the
healing of a man at the Temple's Beautiful Gate, he again addresses Israel,
calls for repentance, and identifies Jesus with the promised prophet and the
servant. When the authorities arrest Peter and John, Acts represents their
boldness as the work of the Spirit in men without elite schooling. Their
confession---that salvation is in Jesus and that they cannot stop speaking of
what they have seen and heard---defines apostolic witness as obedience to God
under pressure.

Peter's authority also appears in severe scenes. He confronts Ananias and
Sapphira over deceit concerning a voluntary gift; both fall dead. He condemns
Simon of Samaria's attempt to acquire spiritual power with money. These
episodes have contributed to Christian language about simony and ecclesial
discipline. They should not be turned into proof that Peter possessed or used
civil coercive power. In Acts, judgment discloses the holiness and freedom of
God's gift; Peter neither executes a sentence nor sells what Simon seeks.

Acts reports healings through Peter and even people placing the sick where his
shadow might fall. It also narrates two imprisonments: release with the
apostles after a nocturnal opening of the prison, and Peter's deliverance by
an angel when Herod Agrippa I has killed James the brother of John. The second
story is told with almost comic concreteness---Peter initially thinks he sees
a vision, a servant named Rhoda leaves him knocking, and the gathered church
does not believe her. Historical reading need not invent a naturalistic
mechanism. Hagiographic reading should notice that deliverance does not make
the apostles invulnerable: James has died, and Peter himself will not always
escape.

\subsection{From Jerusalem to Samaria and the coast}

After Philip evangelizes Samaria, Peter and John come from Jerusalem, pray
for the baptized, and lay hands on them so that they receive the Holy Spirit.
The scene maintains communion between the Jerusalem apostles and a people
long divided from Judea. Peter then travels through Lydda, where Acts reports
the healing of Aeneas, and Joppa, where Tabitha is raised. At Joppa he stays
with Simon, a tanner, a detail suggestive of social and purity boundaries even
before the narrative's decisive vision.

In Acts 10 Peter sees a sheet containing animals and hears the command to kill
and eat. He objects that he has never eaten anything profane or unclean; the
voice answers that what God has cleansed must not be called profane. At that
moment messengers arrive from Cornelius, a Roman centurion described as
devout. Peter enters the Gentile house, announces that God shows no
partiality, and sees the Holy Spirit fall before baptism. The sequence matters:
God acts, and Peter recognizes rather than manufactures the inclusion.

When criticized in Jerusalem for eating with uncircumcised men, Peter retells
the whole event. His defense is not that Israel's Scriptures or holiness no
longer matter, but that he could not hinder God from giving Gentiles the same
gift. Acts 15 later places him at the Jerusalem deliberation, recalling that
God chose Gentiles to hear the gospel through him and arguing against placing
upon them a yoke the Jewish believers and their ancestors could not bear.
James then formulates the judgment and scriptural rationale. Peter speaks
decisively, but he does not act without communal discernment.

\subsection{The historical problem of sequence}

Acts supplies an orderly movement: Cornelius, Jerusalem's acceptance, further
Gentile mission, and the council. Galatians supplies a more abrasive window.
Paul says that Peter already ate with Gentiles at Antioch and then withdrew
when people came from James, fearing the circumcision party. It is tempting
to place the incident either before Acts 10, so that Peter has not yet learned
the lesson, or after Acts 15, so that he culpably violates an agreed policy.
Neither solution can be demonstrated. The correspondence between Paul's
visits in Galatians and Luke's visits in Acts is itself debated.

The safer conclusion is also the more human one. Peter was a major agent in
the Gentile opening and nevertheless found its embodied social consequences
difficult under pressure. Insight, institutional decision, and stable table
practice did not arrive everywhere at once. Acts and Galatians agree that
Gentile inclusion became a defining apostolic question; they let different
dimensions of Peter's role remain visible.

\evidenceaxis{Peter exercised conspicuous leadership in the Jerusalem church
and became a pivotal witness to Gentile inclusion.}{Acts 1--15, read beside
Gal 1--2 and 1 Cor 15.}{Luke's speeches and sequence have a literary shape;
the exact correlation of Acts with Paul's chronology remains contested.}
